
\AddSymbolsD
\pagecolor{Goldenrod!20!Gray}
\color{white}
\quad\quad \lettrine[,lraise=0.2,lhang=0.5]{\textbf{伊}}{藤}佐一为了躲雨，在小路上小跑着。

茶舍的轮廓逐渐显现；一旁的树木威严地守在两边——那是一种罕见的红杉树。

伊藤佐一终于走进了茶舍。他远远地便看到了教授标志性的二郎腿。教授也看到他了，向他招了招手。

他绕过无数茶桌，终于来到教授身旁。铃木和子坐在教授对面，既没有看着教授，也没有看着伊藤佐一；这让他有点失落。

\twkkai{坐吧;}教授向伊藤佐一招了招手。伊藤佐一便坐在了教授身旁。他揣摩着铃木和子的内心世界。

实际上，在走进茶舍的一瞬，伊藤佐一便留意到她满眼的心不在焉；而且她瘦了许多、也憔悴了许多——她的双颊已经有了他远远就能看到的凹陷。

待到他结束他的胡思乱想之后，他才留意到茶舍原来那么嘈杂。

而自从坐下后，他的手便无处安放——一切都是那么地不对劲。

外面下着大雨；而人群依旧吵闹。伊藤佐一的心再也静不下来了。

可是这时，铃木和子的头轻轻地抵在了伊藤佐一的肩上——是的，不顾教授的在场，抵在了他的肩上。伊藤佐一的头动都不敢动；而教授则突然开始宣讲起什么。

\twkkai{
好了，所有关于我噩梦的一切，都发生在同一个舞台——主角就是你们所看到的\(f(\tau)\)和\(\overline{g(\tau)}\)：
\[
(4\pi)^{-s}\Gamma(s)
\frac{L(s,f\otimes g)}
{\zeta(2s-2k+2)}
=
\int_{\Gamma\backslash\mathbb H}
f(\tau)\overline{g(\tau)}
\,y^k
E(\tau,s-k+1)
\frac{dx\,dy}{y^2}.
\]
}

伊藤佐一仿佛又看见了一条鲤鱼；而挨在他肩上的、困意满满的铃木和子，也早已看见了那条长长的鲤鱼。她心中暖暖的，一种陌生的熟悉感涌了上来——不是对公式的熟悉；是对某种感情。

伊藤佐一坐在一旁，静静地看着那条鲤鱼游曳——但这次，他的注意力在教授身上；他仿佛忘了肩膀上的重量。他知道，这是一种幸福。

教授则压根没有留意，继续说着。

\twkkai{
这说明
\[
\text{解析积分}
\quad\Longleftrightarrow\quad
\text{$L$-函数}.
\]
}

听到这里，伊藤佐一意识到什么。

\twkkai{是啊——我怎么没想到这个等价关系呢？当时应该对她讲才对。}

伊藤佐一轻轻地跺着脚，仿佛带着一丝悔意。他依旧不敢动，但他心中多了一个危险的想法——他想用手轻轻地搭在铃木和子的肩上。

窗外吹起了一阵风，扬起了铃木和子的刘海。她再次下意识地挡了挡；她的动作吓跑了伊藤佐一的想法。

\twkkai{好了——那现在就可以讲椭圆曲线了。}

教授顺手写下几个大字。

\begin{center}
\textbf{\twkkai{
椭圆曲线与 $j$-不变量}}
\end{center}

\twkkai{这里的\(j\)-不变量，就是函数\(j(\tau)\)——那么，接下来为了方便，还是写成“\(j\)-不变量”；}

\twkkai{
$j$-不变量是从基本域\(\Gamma\backslash\mathbb H\)到整个复平面\(\mathbb C\)的映射，具体就是这样：
\[
j:\Gamma\backslash\mathbb H\to\mathbb C.
\]
它的定义式是：
\[
\boxed{
j(\tau)
=
\frac{E_4(\tau)^3}{\Delta(\tau)}},\qquad \Delta(\tau)
=
\frac{E_4(\tau)^3-E_6(\tau)^2}{1728}.
\]
铃木同学——我之前应该单独跟你讲过；不知道伊藤同学有没有给你温习。}

铃木和子轻轻点了点头。但一旁的伊藤佐一知道，他没有和她讲过这个被圈起来的定义式；他不知道为什么她要选择撒谎。

\twkkai{
于是——
\(
j(\tau)
\)
完全分类复椭圆曲线。我也和铃木同学讲过——\(j(\tau)\)其实就是椭圆曲线的身份证。}\index{jbubianliang@$j$-不变量 $j(\tau)$!完全分类复椭圆曲线}

伊藤佐一又开始羡慕起铃木和子了——教授到底给她单独分享过多少数学世界里的秘密啊。他开始觉得自己所在的这个现实世界变得更无趣了——无法理解，无法探秘，也仿佛再也无法生存。

但是，又幸好——他还有铃木和子。

\twkkai{很好，很好——接下来便是我的噩梦了；}教授突然说。

伊藤佐一不敢吭声。不仅是因为铃木和子几乎已经睡着；更是因为教授语气的沉重——仿佛自己的任何提问都会显得过于轻浮。

\twkkai{这个噩梦就是——为什么 $e^{\pi\sqrt{163}}$ 几乎是整数？哎，没成想——拉马努金的女神托的梦，反倒变成了我怎么都摆脱不了的梦魇。}

\twkkai{
让我详细地说说吧。如果把自变量从\(\tau\)变为\(e^{2\pi i\tau}\)的话，就会导致刚才这个神奇效应。实际上，这个指数化并不是什么很特殊的操作，只是傅里叶分析的一种基本工具而已。}

教授心不在焉地说着，似乎连笔也挪不动了——伊藤佐一能感受到。茶舍里突然又多了几处笑声。

\twkkai{该死的163——好吧，把163代入得到
\(
q
=
-e^{-\pi\sqrt{163}}.
\)
那么，$j$-函数的傅里叶展开就很自然地得到了：
\[
j(\tau)
=
q^{-1}
+744
+196884q
+21493760q^2
+\cdots.
\]
代入\(q\)得到：
\[
j(\tau_{163})
=
-e^{\pi\sqrt{163}}
+744
-196884e^{-\pi\sqrt{163}}
+\cdots.
\]}

\twkkai{
而另一方面：
\[
j(\tau)=-640320^3.
\]
这个我也和铃木同学讲过的。对吧，铃木同学？}

教授终于注意到沉睡中的铃木和子。茶舍中的笑声越来越嘈杂了。

\twkkai{铃木同学？好吧，看来她是比我还要累了。不过没关系，她如果忘了，也挺好；}教授近乎疯子般喃喃自语着。

伊藤佐一不知道教授为什么说“她忘了更好”这句话；他无法理解。他终于在这时意识到，铃木和子似乎处于很危险的状态。

茶舍的吵闹达到了顶峰。教授突然站起身，去前台拿了一壶茶，斟给伊藤佐一一小杯。伊藤佐一连忙用手指叩了桌面几下。

但教授似乎彻底忘记了伊藤佐一——以及铃木和子——的存在，数起了茶来。

1，2，3，……

7，11，19，……

43，67，……

163……

茶滴停下来了。而他又想起自己还没有兑现的诺言。他无奈地放下了茶壶，看向伊藤佐一。

\twkkai{那……伊藤同学，我和你再讲一次吧；从此以后我就封笔不讲了。刚才\(j(\tau_{163})=-640320^3\)这个现象，是因为163这个判别式导致虚二次域类数为1，所以\(j(\tau_{163})\)是一个整数。如果要硬挖，那就是163为导致这个整数现象的最大的数——即最大的Heegner数。
于是：
\[
e^{\pi\sqrt{163}}
=
640320^3+744+\varepsilon,
\]
其中\(
|\varepsilon|\ll1.
\)
因此：
\(
e^{\pi\sqrt{163}}
\approx
262537412640768744.
\)
这就是著名的近整数现象……啊，真是让人头皮发麻……}\index{jinzhengshuxianxiang@近整数现象 $e^{\pi\sqrt{163}}\approx 640320^3+744$}

教授看着那串数字——262537412640768744——全身打了个寒颤。伊藤佐一突然是那么为教授感到心疼；但教授依旧万般严肃。

\twkkai{好——让我总结一下。同一个空间
\(
\Gamma\backslash\mathbb H
\)
同时承载了
\begin{itemize}
    \item Rankin--Selberg 积分；
    \item 椭圆曲线；
    \item $j$-不变量；
    \item 椭圆曲线 $L$-函数。
\end{itemize}
}\index{rankin@Rankin-Selberg卷积!Rankin-Selberg 积分}


教授的笔顿了顿。他点的那四个点，大小完全不一样；这让此刻的他在伊藤佐一眼中显得有点滑稽。

教授的笔仿佛对伊藤佐一的想法感到不忿气，继续在纸上飞舞。

\twkkai{
那么，整体链条就是这样：
\[
\tau
\longmapsto
E_\tau
\longmapsto
j(\tau).
\]
另一方面：
\[
f, g, E(\tau, s) \longmapsto L(s, f \otimes g),
\]
给出 $L$-函数的解析表示。这就是模函数（注意，不是模形式）\(j(\tau)\)和Heegner数 $163$ 的传奇。
}

伊藤佐一震惊于教授竟还有精力提醒他模函数和模形式区别的事实；他更心疼教授了。教授终于放下了笔，又擦了擦眼镜。接着，像在草地上拍大腿那样——这次，他拍了拍两边衣袖，然后慢慢站起了身，伸了个懒腰，打了个迟到了许久的哈欠。

\twkkai{哎……是163的传奇——也是我该退场的时候了……；}教授突然说出了这一句话。

\twkkai{退场？}伊藤佐一沉思着，看着纸上的\(\Gamma\backslash\mathbb H\)发呆，又想起了自己向铃木和子授课的温馨场景。

他没想到，教授向自己授课这一同样温馨的场景终将告一段落。他终于还是忍不住问出了他的问题。但他觉得这值得一问。

\twkkai{退场？}他问的比他想的要小声得多。

\twkkai{是的；从此之后，我要消失一段时间了。暑假也快来了，学生放假了；至于我，只能回去那个地方了。}

教授伸着懒腰的双手突然多了两根手指，往茶舍的茅草天花顶部指了指。接着，他拿起笔，就要走。伊藤佐一想站起身，但他站不起来。

不过几秒，教授便消失于门口。茶舍逐渐变得安静了。

伊藤佐一后悔自己没有组织好语言，问出一切。不过，铃木和子依旧挨在他的肩上——他还不敢做什么。他害怕这个时刻，同时又想珍惜这个时刻。

……

终于，铃木和子醒了。她捂着头，揉着眼睛，面露疑惑之色。

\twkkai{嗯……？我这是在哪？你是谁？}

\twkkai{我……我是伊藤佐一啊……；}伊藤佐一惶恐地看着她，发现自己刚才想珍惜的那个所谓“幸福”的时刻，原来是一个漏洞。

\twkkai{我不认识你。你为什么会和我在一起？}

伊藤佐一无助地看着那些嬉笑着品茶的人群。不，他们只是一群鲁莽的用茶水灌进自己的胃里的异类而已——无异于酒鬼。酒鬼怎么会同情他如今的处境呢？

\twkkai{你是我的同学吗？}

伊藤佐一听到“同学”二字，突然倍感欣慰——至少她不是什么都不知道。他连忙充满希望地点了点头。

\twkkai{你和我离开这里吧。我有点恶心。我想散散步。}

铃木和子尝试站起，但没有伊藤佐一的搀扶显然不可能。

他们离开了茶舍。

他们走了很远很远。

一片树林在向他们招手。



